\documentclass[12pt]{exam}
\usepackage{libertine}
\usepackage[utf8]{inputenc}
\usepackage[margin=1in]{geometry}
\usepackage{amsmath,amssymb,amsthm}
\usepackage{multicol}
\usepackage[shortlabels]{enumitem}
\usepackage{siunitx}
\usepackage{cancel}
\usepackage{graphicx}
\usepackage{pgfplots}
\usepackage{listings}
\usepackage{tikz}
\usepackage{booktabs}
\usepackage{indentfirst}


\pgfplotsset{width=10cm,compat=1.9}
\usepgfplotslibrary{external}
% Uncomment if your compiler has shell-escape enabled:
% \tikzexternalize

\newcommand{\class}{Research in CompMath} 
\newcommand{\examnum}{Homework 3} 
\newcommand{\examdate}{10/10/2026}

\begin{document}
\pagestyle{plain}
\thispagestyle{empty}

\noindent
\begin{tabular*}{\textwidth}{l @{\extracolsep{\fill}} r @{\extracolsep{6pt}} l}
\textbf{\class} & \textbf{Name:} & \textit{Mary Kim}\\ 
\textbf{\examdate} &&\\
\end{tabular*}\\
\rule[2ex]{\textwidth}{2pt}
% -------

4.1, 4.3, 4.5, 4.7, 4.9, 4.13, 4.15, 4.17 

\bf 4.1 Show by induction that $2^n < n!$ for all $n \ge 4$.
\normalfont
\\For the base statement, n = 4, the statement that $2^4 < (4)!$ holds true as it boils down to $16 < 24$. 
\\An inductive proof for this could be shown below
\\$2^{n+1} < (n+1)!$ which can be rewritten as $2(2^n) < (n+1)!$
\\Since $n > 4$, then $n + 1 > 5$, which means that $n + 1 > 2$ also hold true, multiplying $n!$ results in: $2n! < (n+1)n!$
\\If you multiply the induction hypothesis ($2^n < n!$) by 2, you get: 
\\$2(2^n) < 2n!$
\\Then combining the above inequalities yields:
\\$2^{n+1} < 2n! < (n+1)!$
\\This proves that $2^n < n!$ holds for n+1. \hfill $\square$

\bf 4.3 Show by induction that $f_0 + ... + f_n = f_{n+2} - 1$ in the Fibonacci sequence. 
\normalfont
\\First, proving the base case: 
\\For $n=0$, $f_n = 1$ so $f_{n+2} - 1 = 2-1 = 1$ This proves the base case. 
\\Now, assuming that n is true:
\\$f_n = f_{n+1} - 1$
\\Finally, the inductive step:
\\ 
\\so $f_{n+1} = f_{n+2} - 1$

\bf 4.5 Consider an alternative Fibonacci sequence, starting with 0, 2. Can you prove the analogue of theorem 25? Generalize the result as far as you can. 

\bf 4.7 Show by induction that a finite set with n elements has exactly $2^n$ many subsets. 

\bf 4.9 Prove that $\frac{1}{1*2} + \frac{1}{2*3} + \frac{1}{3*4} + ... + \frac{n}{n*n+1} = \frac{n}{n+1}$.

\bf 4.13 True or false: every flying polka-dotted elephant in this room is smoking a cigar. 

\bf Criticize the following ``proof.'' 

\textit{Claim.} $\sum_{i=1}^{\infty} \frac{1}{i} < \infty$. 

\textit{Proof.} Let $S(n) = \sum_{i=1}^{n} \frac{1}{i}$, and consider the statement asserting $S(n) < \infty$. This statement is true for $n = 1$, since $\frac{1}{1} = 1 < \infty$. And if it is true for $n$, then it is true for $n + 1$, since the next sum $S(n + 1) = \sum_{i=1}^{n+1} \frac{1}{i}$ is equal to the previous sum $\sum_{i=1}^{n} \frac{1}{i}$ plus the next term $\frac{1}{n+1}$, or in other words, $S(n + 1) = S(n) + \frac{1}{n+1}$, and this is the sum of two finite numbers. So $S(1)$ is true, and $S(n)$ implies $S(n + 1)$. So we have therefore proved $\sum_{i=1}^{\infty} \frac{1}{i} < \infty$ by induction. \hfill $\square$

\bf 4.17 Prove that the least-number principle, the common induction principle, and the strong induction principle are all equivalent. Assuming any one of them, you can prove the others. 
\end{document}